\chapter{Постановка задачи и анализ предметной области}
\label{cha:analysis}

\section{Описание предметной области}

Предметная область---управление учебным расписанием высшего учебного заведения. Система оперирует следующими сущностями:

\begin{itemize}
\item \textbf{Подгруппа} (Subgroup)---академическая группа студентов;
\item \textbf{Преподаватель} (Teacher)---сотрудник кафедры, проводящий занятия;
\item \textbf{Аудитория} (Location)---место проведения занятия;
\item \textbf{Дисциплина} (Subject)---учебный предмет;
\item \textbf{Расписание} (Timetable)---семестровый период с указанием дат и типа недели (числитель/знаменатель);
\item \textbf{Занятие} (Lesson)---проведение дисциплины в определённое время, с привязкой к преподавателю, аудитории и подгруппам.
\end{itemize}

\section{Функциональные требования}

На основании анализа предметной области сформулированы следующие функциональные требования:

\begin{itemize}
\item CRUD-операции для всех справочных сущностей;
\item CRUD-операции для занятий с гибким связыванием преподавателей и аудиторий;
\item импорт расписания из ZIP-архива, содержащего XLSX-файлы определённого формата;
\item экспорт расписания в формате iCalendar (.ics) для интеграции с календарями;
\item визуализация расписания в виде недельной сетки;
\item разграничение доступа по ролям.
\end{itemize}

\chapter{Используемые инструменты и технологии}
\label{cha:tools}

\section{Языки программирования и среды разработки}

\textbf{Серверная часть---Go.} Выбор обусловлен высокой производительностью, встроенной поддержкой конкурентности, простым синтаксисом и широкими возможностями для создания HTTP-сервисов~\cite{go:doc}.

\textbf{Десктопный клиент---C++20 с Qt~6.} Выбор Qt~6 обусловлен кроссплатформенностью, богатым набором виджетов и встроенной поддержкой сетевых запросов~\cite{qt:doc}.

\section{Библиотеки и API}

Серверная часть использует следующие основные библиотеки:
\begin{itemize}
\item \Code{pgx/v5}---PostgreSQL драйвер с пулом соединений;
\item \Code{oapi-codegen}---генерация Go-кода из OpenAPI~\cite{oapi:codegen};
\item \Code{excelize/v2}---чтение и парсинг XLSX-файлов;
\item \Code{golang-ical}---генерация iCalendar;
\item \Code{goose/v3}---миграции базы данных~\cite{goose:migrate};
\item \Code{sqlc}---генерация type-safe Go-кода из SQL~\cite{sqlc:dev}.
\end{itemize}

Десктопный клиент на C++/Qt~6 использует:
\begin{itemize}
\item Qt~6 (Core, Gui, Widgets)---построение GUI и сетевое взаимодействие;
\item OpenAPI Generator (cpp-qt-client)---генерация C++~Qt-клиента из OpenAPI~\cite{openapi:generator}.
\end{itemize}

\section{Инструменты развёртывания}

\begin{itemize}
\item \textbf{Docker}---контейнеризация сервисов~\cite{docker:doc};
\item \textbf{Docker Compose}---оркестрация многоконтейнерного приложения;
\item \textbf{PostgreSQL}---реляционная база данных~\cite{postgres:doc};
\item \textbf{Nginx}---обратный прокси и раздача статических файлов.
\end{itemize}
